\scanpage{7}
\[
\begin{aligned}
(11)\qquad\qquad & \underline{\operatorname{Lie}}(P)\handwritten{\simeq}H^1(X,\underline{O}_X)\ .
\end{aligned}
\]
As $R^{\mathord{\rlap{$1$}\handwritten{\unsure{4|1}}}}f_{\handwritten{\unsure{*|k}}}(\underline{O}_X)$\unsure{;} is coherent and concentrated at the origin, it follows that its group of sections $H^1(X,\underline{O}_X)$ is of finite dimension. In fact, it seems, with a little care, that \deleted{\unsure{get|gei}} replacing if necessary $X$ by a \unsure{snitable|suitable} model \unsure{dmminating|dominating} it, we may assume that there exists an \tikz[baseline=(s.base)]{\node[inner sep=1pt](s){ideal};\node[anchor=south west,inner sep=1pt](a) at ([xshift=-4em,yshift=7pt]s.north west){\inserted{invertible}};\draw ([xshift=-4pt,yshift=-3pt]s.south west) to[out=90,in=-90] ([xshift=-4pt,yshift=-2pt]s.north west);\draw ([xshift=-3pt,yshift=5pt]a.south west) to[out=210,in=180] ([xshift=-4pt,yshift=-2pt]s.north west) to[out=0,in=180] ([xshift=1pt,yshift=-2pt]a.south east) to[out=0,in=-90] ([xshift=3pt,yshift=4pt]a.south east);} $J$ on $X$, \deleted{having} defining a subscheme $Z$ of $X$ having same support as $X_o$, and such that $\mathord{\tikz[baseline=(s.base)]{\node[inner sep=0pt](s){$J/J^2$};\node[anchor=south east,inner sep=1pt](a) at ([xshift=0pt,yshift=8pt]s.north west){$\inserted{\underline{L}=}$};\draw ([xshift=-4pt,yshift=-3pt]s.south west) to[out=90,in=0] ([xshift=-1pt,yshift=-2pt]a.south west) to[out=180,in=-90] ([xshift=-4pt,yshift=4pt]a.south west);\draw ([xshift=-4pt,yshift=-3pt]s.south west)--([xshift=9pt,yshift=16pt]s.north west);}}$ be an \uline{ample sheaf on} $Z$. This will imply that, replacing the system $(X_n)_{\handwritten{n}}$ by the system $(Z_n=V(J^{n+1}))_{\handwritten{n}}$, which defines \deleted{the} \deleted{\unsure{same}} an isomorphic pro-\rlap{\deleted{group}}\handwritten{object} of $\unsure{P|\mathbb{P}}$icard -schemes, we obtain an inverse system of group-schemes $P'_n=\underline{\operatorname{Pic}}_{Z_n/k}$ where for large $n$ the transition morphisms \deleted{are \unsure{h|b}} $P'_{n+1}\quad P'_n$ are \uline{isomorphisms} (because $H^1(Z_o,L^{\mathord{\tikz[baseline=(s.base)]{\node[inner sep=0pt](s){$\scriptstyle\mathrm{x}$};\draw[line width=.4pt] ([xshift=.7pt,yshift=1.2pt]s.center) ellipse[x radius=1.8pt,y radius=3.1pt];}}n})=H^2(Z_o,\underline{L}^{\mathord{\tikz[baseline=(s.base)]{\node[inner sep=0pt](s){$\scriptstyle\mathrm{x}$};\draw[line width=.4pt] ([xshift=.7pt,yshift=1.2pt]s.center) ellipse[x radius=1.8pt,y radius=3.1pt];}}n})=0$). \unsure{Thhs|Thus} it seems that, under the assumptions made \tikz[baseline=(s.base)]{\node[inner sep=1pt](s){,};\node[anchor=south west,inner sep=1pt](a) at ([xshift=-3em,yshift=9pt]s.north west){\inserted{\handwritten{from the start}}};\draw ([xshift=-2pt,yshift=-4pt]s.south west) to[out=90,in=180] ([xshift=1pt,yshift=-2pt]a.south east) to[out=0,in=-90] ([xshift=3pt,yshift=4pt]a.south east);} one can prove that the pro-group $(P_n)$ is essentially constant, hence its inverse limit $P$ is in fact a \uline{group scheme locally of finite type over} $k$. This is certainly so, at least, if $\dim A=2$, $A$ normal, as follows from the negative definiteness of the intersection matrix of the components of $X_o$.

\unsure{This|Thus|Thhs} $P$ seems to be definable in quite a reasonable way, and one still has to investigate the morphism (5) (or more accurately, \tikz[baseline=-.3ex,x=1em,y=1em]{\fill (.12,-.15) circle (.3pt);\fill (.24,0) circle (.3pt);} the morphisme $\underline{D}\handwritten{\longrightarrow}P$ defined by descent from the \rlap{x}case where $k$ is separably \tikz[baseline=(c.base)]{\node[anchor=base west,inner sep=0pt](c){closed;};\node[anchor=base west,inner sep=0pt](b) at ([xshift=.4em]c.base east){b\rlap{y}ut};\node[anchor=base west,inner sep=0pt](t) at ([xshift=.4em]b.base east){then};\node[anchor=base west,inner sep=0pt](w) at ([xshift=.8em]t.base east){we};\draw[gray,line width=.25pt] ([xshift=-1.8em,yshift=.8em]c.west)--([xshift=-.4em,yshift=-.8em]c.west) to[out=-5,in=185] ([yshift=-.6em]b.south) to[out=5,in=180] ([xshift=-.2em]t.west);\draw[line width=.35pt] ([xshift=-.2em,yshift=.1em]t.east) to[out=-60,in=165] ([xshift=-.1em,yshift=-.1em]w.west);} may as well assume $k$ separably closed \rlap{g}from the start\handwritten{,}\tikz[baseline=(s.base)]{\node[inner sep=1pt](s){.};\node[anchor=south west,inner sep=1pt](a) at ([xshift=-2em,yshift=7pt]s.north west){\inserted{\handwritten{which we will do}}};\draw ([xshift=-4pt,yshift=-3pt]s.south west) to[out=90,in=-90] ([xshift=-4pt,yshift=-2pt]s.north west);\draw ([xshift=-3pt,yshift=5pt]a.south west) to[out=210,in=180] ([xshift=-4pt,yshift=-2pt]s.north west) to[out=0,in=180] ([xshift=1pt,yshift=-2pt]a.south east) to[out=0,in=-90] ([xshift=3pt,yshift=4pt]a.south east);}). We would be particularly happy if this were a closed immersion, or what amounts to the same, if the morphism\deleted{\unsure{s|e}}
\[
\begin{aligned}
\handwritten{(12)}\qquad\qquad & D\mathrel{\tikz[baseline=-.3ex,x=1em,y=1em]{\draw[line width=.5pt] (0,0)--(1.5,0)--(1.32,.10) (1.5,0)--(1.32,-.13);\fill (1.22,-.15) circle (.35pt);}}\operatorname{NS}(X_o)
\end{aligned}
\]
is injective, were $\operatorname{NS}(X_o)=\underline{\operatorname{NS}}_{X_o\mathord{\tikz[baseline=-.25ex,x=1em,y=1em]{\draw[line width=.7pt] (0,-.05)--(.32,.72);\fill (.13,.45) ellipse[x radius=.06em,y radius=.08em];\fill (.14,.18) ellipse[x radius=.06em,y radius=.08em];}}k}(k)$, $\underline{\operatorname{NS}}=\underline{\operatorname{Pic}}/\underline{\operatorname{Pic}}^o$. (The equivalence comes from the fact that $P=\underline{\operatorname{Pic}}_{Z_n/k}$ for some \handwritten{l}arg\handwritten{e} $n$, and that $\deleted{\operatorname{Pic}}$
% archive sheet number: 6
